\section{Произношение гласных звуков}

\begin{longtblr}{
    colspec = { X[1,c,t] X[1,c,t] X[1,c,t] X[4,l,t] },
    rowhead=1,
    row{1} = {font=\bfseries},
}
\toprule
Румынский звук & Русский аналог & Написание & Комментарий \\
\midrule
\textipa{[I]}
& и
& i
&
\\
\textipa{[e]}
& { э \\ (закрытое) }
& e
& { В конечном открытом слоге (оканчивающемся на гласную) произносится закрыто \textipa{[e]}
  (уголки губ растягиваются в стороны, как при улыбке). \\
  В многосложных словах закрытое произношение \textipa{[e]} распространяется,
  как правило, и на предыдущие открытые слоги. }
\\
\textipa{[E]}
& э
& e
& В остальных случаях, которые не описаны в произношении звука \textipa{[e]},
  произносится \textipa{[E]}, близкое к русскому звуку \textipa{[\textnormal{э}]}.
\\
\textipa{[u]}
& у
& u
&
\\
\textipa{[o]}
& о
& o
&
\\
\textipa{[1]}
& ы
& î, â
&
\\
\textipa{[a]}
& а
& a
&
\\
\textipa{[@]}
& --
& ă
& Получить этот звук можно произнося румынский \textipa{[o]},
  сняв при этом его огубленность (рекомендуется приложить к губам карандаш).
\\
\bottomrule
\end{longtblr}

Примеры:
\begin{description}
    \item[\textipa{[I]}:] vin \textipa{[vIn]}, spirit \textipa{["spIrIt]}, finit \textipa{[fI"nIt]}.
    \item[\textipa{[e]}:] vine \textipa{["vIne]}, bine \textipa{["bIne]}, primite \textipa{[prI"mIte]};
        sete \textipa{["sete]}, perete \textipa{[pe"rete]}, repede \textipa{["repede]}.
    \item[\textipa{[E]}:] drept \textipa{[drEpt]}, metru \textipa{["mEtru]}, eveniment \textipa{[EvEnI"mEnt]}.
    \item[\textipa{[u]}:] bun \textipa{[bun]}, drum \textipa{[drum]}, grup \textipa{[grup]}.
    \item[\textipa{[o]}:] bob \textipa{[bob]}, nostru \textipa{["nostru]}, popor \textipa{[po"por]}.
    \item[\textipa{[1]}:] român \textipa{[ro"m1n]}, România \textipa{[rom1"nIa]},
        în \textipa{[1n]}, întreb \textipa{[1n"trEb]}.
    \item[\textipa{[a]}:] amar \textipa{[a"mar]}, armata \textipa{[ar"mata]}, strada \textipa{["strada]}.
    \item[\textipa{[@]}:] ăl \textipa{[@l]}, ăst \textipa{[@st]}; măr \textipa{[m@r]}, păr \textipa{[p@r]};
        apă \textipa{["ap@]}, masă \textipa{["mas@]}, bărbat \textipa{[b@r"bat]}.
\end{description}
